On compte le nombre d'arbres de deux forêts où de nouveaux arbres sont plantés chaque année.

\subsection*{Partie A :\quad Première forêt}

Dans cette forêt, au 1\up{er} janvier 2010, on dénombrait \np{1200}~arbres.

On suppose que le nombre d'arbres augmente de \np{100} par an.

On modélise cette évolution par une suite \mbox{$(u_n)$}, où $u_n$ représente le nombre d'arbres dans cette forêt au 1\up{er}~janvier de l'année \mbox{$2010+n$} pour tout entier naturel $n$.

Ainsi, on a \mbox{$u_0=\np{1200}$}
\begin{enumerate}
	\item Calculer $u_1$ et $u_2$.

	Interpréter $u_2$ dans le contexte de l'exercice.
	\item Exprimer \mbox{$u_{n+1}$} en fonction de $u_n$ pour tout entier naturel $n$.

	En déduire la nature de la suite $(u_n)$ et préciser sa raison $r$.

	\item Exprimer $u_n$ en fonction de $n$.

	\item D'après ce modèle, à partir de quelle année la forêt comptera-t-elle plus de \np{2950}~arbres ?
\end{enumerate}

\subsection*{Partie B : \quad Seconde forêt}


Dans cette forêt, au 1\up{er} janvier 2010, on dénombrait \np{1000}~arbres.

On suppose que le nombre d'arbres augmente de 5\,\% par an.

On modélise cette évolution par une suite \mbox{$(v_n)$}, où $v_n$ représente le nombre d'arbres dans cette forêt au 1\up{er}~janvier de l'année \mbox{$2010+n$} pour tout entier naturel $n$.

Ainsi, on a \mbox{$v_0=\np{1000}$}

\begin{enumerate}
	\item Calculer le nombre d'arbres au 1\up{er}~janvier 2011.
	\item Déterminer la nature de la suite \mbox{$(v_n)$} et préciser sa raison $r$.
	\item D'après ce modèle, on estime qu'au 1\up{er}~janvier 2030, il y aura environ \np{2653}~arbres sur la parcelle.

	Quel terme de la suite a-t-on calculé?

	Quel a été le calcul effectué ?
\end{enumerate}

\subsection*{Partie C : \quad Comparaison des évolutions}
Grâce à l'extrait de tableur ci-dessous, on visualise le nombre d'arbres arrondi à l'unité dans les deux forêts.

\begin{center}
	\begin{tblr}{|X[8mm]|*{9}{X[c]|}}\hline
$n$&	0&	1&	2&	3&	\dots&	26&	27&	28&	\textbf{29}\\\hline
$u_n$&	\np{1200}&	\np{1300}&	\np{1400}&	\np{1500}&	\dots&	\np{3800}&	\np{3900}&	\np{4000}&	\textbf{\np{4100}}\\\hline
$v_n$&	\np{1000}&	\np{1050}&	\np{1103}&	\np{1158}&	\dots&	\np{3556}&	\np{3733}&	\np{3920}&\textbf{	\np{4116}}\\\hline
\end{tblr}

\end{center}
Interpréter dans le contexte de l'exercice les résultats de la dernière colonne.
